\subsection{Complete particle snapshots: phase-space and energy distributions}
\label{sec:complete-particle-analysis}

The GPU particle runners can save the complete marker population together
with the gravitational field at selected times. Each marker then has an orbital state and a potential evaluated at the
same time, so the complete mass-weighted orbital distribution can be measured.

Use \code{analysis/scripts/phase_space_snapshot_analysis.py} to calculate the
phase-space distributions.

Use \code{analysis/scripts/shifted_energy_snapshot_analysis.py} to calculate the
energy distributions. Both tools run on a CPU and require NumPy.

A complete group contains \code{particles.bin}, \code{field.json}, and
\code{manifest.json}. The supported native format is the version-6 GPU
particle format: a 112-byte header followed by 81 bytes per marker. The
header stores the accepted time, step, counters, and event clocks. Each
marker supplies mass $m_p$, physical radius $r_p$, peculiar radial
velocity $w_p$, and $\eta_p=L_p^2/2$, where $L_p$ is specific angular
momentum. Internal energy is stored separately. The supported population
has active positive-mass markers at positive radii within the outer field
face. The reader checks the exact file size, complete SHA256 hashes,
count and clock agreement, finite values, and unique marker identifiers.

The field file supplies physical cell faces, source mass at those faces,
constant source line density within each cell, the potential $\Phi$ and
its radial derivative, and the prescribed expansion rate $H$. For the
supported positive source with homogeneous subtraction disabled, the
reader verifies the enclosed-mass recurrence, the hidden central mass,
and the complete deposited-mass budget. It independently integrates this
source to check the saved potential and force. The potential is normalized
to zero at infinity. Static total-mass snapshots and the expanding
runners' directly supplied positive contrast-mass snapshots use this
format. The shifted-energy calculation also requires zero central point
mass, giving a finite central potential.

At a fixed time the physical radial velocity is $u=w+Hr$. This velocity
shift has unit Jacobian, so the mass distribution $f$ satisfies
\begin{equation}
 \boxed{\mathrm dM=8\pi^2 f\,\mathrm dr\,\mathrm du\,\mathrm d\eta.}
 \label{eq:df-analysis-measure}
\end{equation}
Integrating over angular momentum and averaging over a radial--velocity
bin gives the discrete projection
\begin{equation}
 \boxed{\left\langle\int f\,\mathrm d\eta\right\rangle_{ij}
 =\frac{\sum_{p\in(i,j)}m_p}{8\pi^2\Delta r_i\Delta u_j}.}
 \label{eq:df-analysis-phase-bin}
\end{equation}
Here $\Delta r_i$ and $\Delta u_j$ are the physical bin widths. A
logarithmic radius axis still uses
$\Delta r_i=10^{x_{i+1}}-10^{x_i}$, where $x_i$ is a logarithmic edge.
The code accumulates the mass and squared mass before forming the
projection and the effective population
$N_{\rm eff}=(\sum m_p)^2/\sum m_p^2$. In
\code{analysis/scripts/phase_space_snapshot_analysis.py:709}, these operations
are shown here, with \code{PHASE_MEASURE} equal to $8\pi^2$.
\codesource{Mass-weighted phase portrait}{\path{analysis/scripts/phase_space_snapshot_analysis.py}, \code{histogram}, lines 709--718}{Divide each marker-mass bin by its physical phase measure and calculate its effective population}
\begin{lstlisting}[language=Python]
    mass = mass.reshape(shape)
    mass2 = mass2.reshape(shape)
    # Use physical radial widths to recover the eta-integrated distribution from mass bins.
    dr = np.diff(10**xedge)
    du = np.diff(uedge)
    dx = np.diff(xedge)
    f_eta = mass / (PHASE_MEASURE * dr[:, None] * du[None, :])
    # Keep the mass density per log-radius interval separately from the distribution per physical dr.
    mass_log_radius = mass / (dx[:, None] * du[None, :])
    neff = np.divide(mass * mass, mass2, out=np.zeros_like(mass), where=mass2 > 0)
\end{lstlisting}

The same markers can be grouped by instantaneous specific orbital energy,
\begin{equation}
 \boxed{E=\frac{(w+Hr)^2}{2}+\frac{\eta}{r^2}+\Phi(r,t).}
 \label{eq:df-analysis-energy}
\end{equation}
The background and halo field can evolve, so this energy changes during
the simulation. The analysis uses the saved field at the marker time.
For a subhalo marker, $E$ describes its center-of-mass motion; the
subhalo's stored internal energy remains a separate quantity.

Counting marker mass in energy bins gives a mass spectrum. To infer a mean
phase-space density, divide that mass by the phase volume available within
each energy bin. Let $\Phi_{\rm c}=\Phi(0,t)$, and let $R_{\rm wall}$ be the outer saved
field face. The spherical phase volume below energy $E$ is
\begin{equation}
 \boxed{\Omega(E)=\frac{16\pi^2}{3}
 \int_0^{R_{\rm wall}}r^2[2(E-\Phi(r,t))]_+^{3/2}\,\mathrm dr,}
 \label{eq:df-analysis-phase-volume}
\end{equation}
where $[x]_+=\max(x,0)$. This follows by integrating the velocity-space
sphere allowed at each physical radius. For energy bin $k$, with edges
$E_k,E_{k+1}$ and marker mass $M_k$, the discrete quantities are
\begin{equation}
 \boxed{\left.\frac{\mathrm dM}{\mathrm dE}\right|_k
 =\frac{M_k}{E_{k+1}-E_k},\qquad
 \langle f\rangle_{E,k}=\frac{M_k}{\Omega(E_{k+1})-\Omega(E_k)}.}
 \label{eq:df-analysis-energy-bin}
\end{equation}
The mean $\langle f\rangle_{E,k}$ applies to the finite-time distribution,
including velocity anisotropy. The logarithmic energy coordinate is
$E-\Phi_{\rm c}>0$; the mass spectrum retains the linear energy-bin width.

The phase-volume integral is divided at each saved radial face and at
the turning radius where $E=\Phi(r,t)$. Gauss--Legendre quadrature uses
24 and 48 nodes per segment; the relative difference of every finite-bin
volume must be below $10^{-6}$. On the final segment, with lower and upper radii $a$ and $b$, the change of
coordinate $r=b-(b-a)(1-z)^2$, $0\le z\le1$, resolves the velocity cutoff
at the upper endpoint. The numerical quadrature and final division are
implemented by the following two operations. Let $r_s(z)$ be the radial
mapping on segment $s$, and let $z_\ell$ and $\omega_\ell$ be the $q$
Gauss--Legendre nodes and weights on $[0,1]$. The discrete integral is
\begin{equation}
 \boxed{\Omega_q(E)=\frac{16\pi^2}{3}
 \sum_s\sum_{\ell=1}^q\omega_\ell r_s^2(z_\ell)r_s'(z_\ell)
 [2(E-\Phi(r_s(z_\ell),t))]_+^{3/2}.}
 \label{eq:df-analysis-quadrature}
\end{equation}
Here $r_s'(z)=\mathrm dr_s/\mathrm dz$. In the first excerpt,
\code{z} and \code{weights} are Gauss--Legendre nodes and weights on
$[0,1]$, \code{value} is $E-\Phi_{\rm c}$, and
\code{self.offset(r)} is $\Phi(r)-\Phi_{\rm c}$.
\codesource{Finite spherical phase volume}{\path{analysis/scripts/shifted_energy_snapshot_analysis.py}, \code{PhaseVolume.integrate}, lines 74--86}{Integrate the allowed velocity-space sphere over saved radial segments}
\begin{lstlisting}[language=Python]
            r = a + (b - a) * z
            jac = np.broadcast_to(b - a, r.shape).copy()
            # The last segment ends at the turning radius (or physical wall).
            r[-1] = b[-1] - (b[-1] - a[-1]) * (1 - z) ** 2
            jac[-1] = 2 * (b[-1] - a[-1]) * (1 - z)
            # Only radii with positive available kinetic energy contribute to the velocity-ball phase volume.
            speed = np.sqrt(2 * np.maximum(value - self.offset(r), 0))
            # Combine the spherical radial measure with the coordinate-map Jacobian and quadrature weights.
            measure = r * r * jac * weights
            cumulative.append(
                (16 * math.pi**2 / 3)
                * float(np.sum(measure * speed**3, dtype=np.longdouble))
            )
\end{lstlisting}
Subtracting cumulative volumes at neighboring energy edges gives the phase
volume of that bin. The code compares both quadrature orders because this
volume sets the normalization of the inferred distribution.
The final division uses \code{mass} for $M_k$ and
\code{delta_volume} for $\Omega(E_{k+1})-\Omega(E_k)$.
\codesource{Energy spectrum and mean distribution}{\path{analysis/scripts/shifted_energy_snapshot_analysis.py}, \code{main}, lines 339--359}{Check finite-bin quadrature and divide marker mass by energy width or phase volume}
\begin{lstlisting}[language=Python]
        coarse, _, _ = volume.integrate(edge, 24)
        omega, _, turn_edge = volume.integrate(edge, 48)
        # A bin volume is the difference of cumulative phase volumes, rather than g at its center times bin width.
        delta_volume = np.diff(omega)
        coarse_volume = np.diff(coarse)
        assert np.all(delta_volume > 0)
        error = np.abs(delta_volume / coarse_volume - 1)
        assert np.max(error) < 1e-6, ("phase quadrature", float(np.max(error)))
        base.require(
            np.all(delta_volume > 0) and np.max(error) < 1e-6,
            "Finite-bin phase-volume quadrature failed.",
        )
        _, g_point, turn_center = volume.integrate(center, 48)
        neff = np.divide(
            mass * mass, squared, out=np.zeros_like(mass), where=squared > 0
        )
        dmass = mass / np.diff(edge)
        # This is the phase-volume mean on the instantaneous finite energy bin.
        # It applies to the sampled distribution during nonequilibrium evolution.
        # Divide the measured bin mass by its actual phase volume to obtain the bin-averaged distribution.
        mean_f = mass / delta_volume
\end{lstlisting}
 The source is the same saved constant-line-mass field
used for the marker energies. Every marker contributes to the cached
histograms; the output also records $N_{\rm eff}$, turning radii, and the
fraction of each energy bin's phase volume inside the first field face.
These quantities identify the population and central continuation used
by a plotted curve.

To generate complete snapshots, add
\code{--phase-space-snapshot-times 0,30,100,300} to a GPU runner with an
initial state at $t=0$ and profile spacing 10. Each requested time must be
a profile time. The selected groups are written below the run's
\code{phase_space/} directory. From the repository root, run
\codesource{Complete-snapshot reduction commands}{\path{analysis/scripts/phase_space_snapshot_analysis.py} and \path{analysis/scripts/shifted_energy_snapshot_analysis.py}, \code{main}}{Read authenticated GPU snapshot groups and write complete phase-space and energy reductions}
\begin{lstlisting}[language=bash]
python analysis/scripts/phase_space_snapshot_analysis.py \
  --snapshots output/my_halo/phase_space \
  --output-directory analysis/my_halo_phase
python analysis/scripts/shifted_energy_snapshot_analysis.py \
  --snapshots output/my_halo/phase_space \
  --output-directory analysis/my_halo_energy
\end{lstlisting}
The first command writes phase-space and linear-energy histograms. The
second writes logarithmic shifted-energy bins and their phase-volume
means. Both write one NPZ file per snapshot, an analysis manifest, and
analytic-check results. The default phase mesh has 360 radial bins and
240 radial-velocity bins; the default shifted-energy mesh has 192 bins.
The command-line bin controls and \code{--edges-from previous_snapshot.npz}
allow common bins to be used in comparisons. Optional expected input,
executable, and metadata SHA256 values enforce the identities chosen by
the reader. The complete files are checked again after reduction.
\code{analysis/README.md} gives a complete run example, array names, plotting
commands, and the supported input conditions.
